% Sample file for AES paper
\documentclass{aes2e}

% Metadata Information
\jyear{2010}
\jmonth{October}
\jvol{1}
\jnum{1}


\begin{document}

% Page heads
\markboth{A1LASTNAME AND A2LASTNAME}{SPECTRAL DELAY FILTERS}


% Title portion
\title{Spectral Delay Filters\thanks{To whom correspondence should be addressed Tel: +1-240-381-2383; Fax: +1-202-508-3799; e-mail: info@schtm.org}}

%Author Info.
\authorgroup{
\author{A1FIRSTNAME A1LASTNAME},
\role{AES Member}
AND \author{A2FIRSTNAME A2LASTNAME},
\role{AES Fellow}
\email{(abc@abc.com)\quad\quad\quad\quad\quad\quad\quad\quad\quad\quad\quad\quad (xyz@abc.com)}
\affil{Universityxyz, City, Country}
}

%Abstract
\abstract{%
This paper discusses the implementation of spectral delay using
filters comprising a cascade of many low-order allpass filters and an
equalizing filter. The spectral delay filters have chirp-like impulse
responses causing a large, frequency-dependent delay that is useful in
audio effects processing. An equalizing filter design and a multirate
technique, which stretches the allpass filters, impulse
response, are introduced.}


\maketitle

%Head 1
\section{INTRODUCTION}
Filtering an audio signal with an allpass filter does not usually have a major effect on the signal's timbre. The allpass filter does not change the frequency content of the signal, but only introduces a phase shift or delay. Audibility of the phase distortion caused by an allpass filter in a sound reproduction system has been a topic of many studies, see, e.g., \cite{DEK1}, \cite{DEK2}. In this paper, we investigate audio effects processing using high-order allpass filters that consist of many cascaded low-order allpass filters. These filters have long chirp-like impulse responses. When audio and music signals are processed with such a filter, remarkable changes are obtained that are similar to the spectral delay effect  \cite{DEK3}, \cite{DEK4}.

\section{CHIRP-LIKE IMPULSE RESPONSES AND GROUP DELAY}
Filtering an audio signal with an allpass filter does not usually have a major effect on the signal's timbre. The allpass filter does not change the frequency content of the signal, but only introduces a phase shift or delay. Audibility of the phase distortion caused by an allpass filter in a sound reproduction system has been a topic of many studies, see, e.g., \cite{DEK1}, \cite{DEK2}. In this paper, we investigate audio effects processing using high-order allpass filters that consist of many cascaded low-order allpass filters. These filters have long chirp-like impulse responses.  These filters have long chirp-like impulse responses.  These filters have long chirp-like impulse responses. When audio and music signals are processed with such a filter, remarkable changes are obtained that are similar to the spectral delay effect  \cite{DEK3}, \cite{DEK4}.

%Head 2
\subsection{Chirp-Like Impulse Responses and Group Delay}
Filtering an audio signal with an allpass filter does not usually have a major effect on the signal's timbre. The allpass filter does not change the frequency content of the signal, but only introduces a phase shift or delay. Audibility of the phase distortion caused by an allpass filter in a sound reproduction system has been a topic of many studies, see, e.g., \cite{DEK1}, \cite{DEK2}. In this paper, we investigate audio effects processing using high-order allpass filters that consist of many cascaded low-order allpass filters. These filters have long chirp-like impulse responses. When audio and music signals are processed with such a filter, remarkable changes are obtained that are similar to the spectral delay effect  \cite{DEK3}, \cite{DEK4}.

%Head 3
\subsubsection{Chirp-Like Impulse Responses and Group Delay}
Filtering an audio signal with an allpass filter does not usually have a major effect on the signal's timbre. The allpass filter does not change the frequency content of the signal, but only introduces a phase shift or delay. Audibility of the phase distortion caused by an allpass filter in a sound reproduction system has been a topic of many studies, see, e.g., \cite{DEK1}, \cite{DEK2}. In this paper, we investigate audio effects processing using high-order allpass filters that consist of many cascaded low-order allpass filters.  When audio and music signals are processed with such a filter, remarkable changes are obtained that are similar to the spectral delay effect  \cite{DEK3}, \cite{DEK4}.
%Equation
\begin{equation}
A(z) = \frac{{a_1  + z^{ - 1} }}{{1 + a_1 z^{ - 1} }},
\end{equation}



Filtering an audio signal with an allpass filter does not usually have a major effect on the signal's timbre. The allpass filter does not change the frequency content of the \nobreak signal, but only introduces a phase shift or delay.\footnote{This point is emphasized by Loewer, see esp. p. (610).} Audibility of the phase distortion caused by an allpass filter in a sound reproduction system has been a topic of many studies, see, e.g., \cite{DEK1}, \cite{DEK2}. In this paper, we investigate audio effects processing using high-order allpass filters that consist of many cascaded low-order allpass filters. These filters have long chirp-like impulse responses. When audio and music signals are processed with such a filter, remarkable changes are obtained that are similar to the spectral delay effect  \cite{DEK3}, \cite{DEK4}.
\begin{equation}
\tau _{\textrm{g,max}}  = \left\{ \begin{array}{l}
 \tau _\textrm{g} (0) = \frac{{1 - a_1 }}{{1 + a_1 }},\textrm{when }a_1  \le 0 \\[4pt]
 \tau _\textrm{g} (\pi ) = \frac{{1 + a_1 }}{{1 - a_1 }},\textrm{when }a_1  > 0. \\
 \end{array} \right.\end{equation}

Filtering an audio signal with an allpass filter does not usually have a major effect on the signal's timbre. The allpass filter does not change the frequency content of the signal, but only introduces a phase shift or delay. Audibility of the phase distortion caused by an allpass filter in a sound reproduction system has been a topic of many studies, see, e.g., \cite{DEK1}, \cite{DEK2}. In this paper, we investigate audio effects processing using high-order allpass filters that consist of many cascaded low-order allpass filters. These filters have long chirp-like impulse responses. When audio and music signals are processed with such a filter, remarkable changes are obtained that are similar to the spectral delay effect  \cite{DEK3}, \cite{DEK4}.
%Paragraph listing
\begin{paralist}
\item{}Filtering an audio signal with an allpass filter does not usually have a major effect on the signal's timbre. The allpass filter does not change the frequency content of the signal, but only introduces a phase shift or delay. 
\item{}Audibility of the phase distortion caused by an allpass filter in a sound reproduction system has been a topic of many studies, see, e.g., \cite{DEK1}, \cite{DEK2}. 
\item{}In this paper, we investigate audio effects processing using high-order allpass filters that consist of many cascaded low-order allpass filters. These filters have long chirp-like impulse responses. 
\item{}When audio and music signals are processed with such a filter, remarkable changes are obtained that are similar to the spectral delay effect  \cite{DEK3}, \cite{DEK4}.
\end{paralist}

Filtering an audio signal with an allpass filter does not usually have a major effect on the signal's timbre. The allpass filter does not change the frequency content of the signal, but only introduces a phase shift or delay. Audibility of the phase distortion caused by an allpass filter in a sound reproduction system has been a topic of many studies, see, e.g., \cite{DEK1}, \cite{DEK2}. In this paper, we investigate audio effects processing using high-order allpass filters that consist of many cascaded low-order allpass filters. These filters have long chirp-like impulse responses. When audio and music signals are processed with such a filter, remarkable changes are obtained that are similar to the spectral delay effect  \cite{DEK3}, \cite{DEK4}.
\begin{arabiclist}
\item{}Green--function determined experimentally and published.
\item{}Black--function determined using similarity searches and published.
\item{}Red--function determined using similarity searches and determined in this study.
\item{}Blue--O-antigen structure unknown. Function determined using similarity searches and proposed in this study.
\end{arabiclist}

%Table
\begin{table}
\tabcolsep8.1pt
\tbl{Active sites and allosteric sites of the GNE MNK enzyme}{%
\begin{tabular}{@{}lccc@{}}\toprule
Excerpt No.& Genre & Spatial Mode & Corrlation\\\colrule
 1 & Pop       & FB   & 94\%\\
 2 & Classical & FB   & 33\%\\
 3 & Jazz      & FF   & 76\%\\
 4 & Arabian   & FF   & 41\%\\
 5 & GNE       & H220 & 45\%\\
 6 & GNE       & H45  & 93\%\\
 7 & MNK       & G416 & 74\%\\
 8 & MNK       & D413 & 72\%\\
 9 & MNK       & R420 & 94\%\\
10 & MNK       & N516 & 91\%\\\botrule
\end{tabular}}
\begin{tabnote}
Note. This table does not include sentence enhancement statutes.  This table does not include sentence enhancement statutes.
\end{tabnote}
\end{table}

%Figure
\begin{figure}
\centering
\includegraphics{aes2e-mouse.eps}
\caption{The spectral delay filter consists of \textit{M} allpass filters and an equalization filter.}
\end{figure}


\begin{figure*}
\centering
\includegraphics[width=23pc]{aes2e-mouse.eps}
\caption{This paper is organized as follows. In Section 1, we discuss the group delay of a cascade of first-order allpass filters and its relation to the chirp-like impulse response of the spectral delay filter. Furthermore, a multirate method to stretch the impulse response of the spectral delay filter is proposed. Section 2 discusses the amplitude envelope of the impulse response and suggests a design method for the equalizing filter. Section 3 presents application examples using the spectral delay filter. Section 4 concludes this paper.}
\end{figure*}
Filtering an audio signal with an allpass filter does not usually have a major effect on the signal's timbre. The allpass filter does not change the frequency content of the signal, but only introduces a phase shift or delay.
\begin{extract}
Filtering an audio signal with an allpass filter does not usually have
a major effect on the signal's timbre. The allpass filter does not
change the frequency content of the signal, but only introduces a
phase shift or delay. Audibility of the phase distortion caused by an
allpass filter in a sound reproduction system has been a topic of many
studies, see, e.g., \cite{DEK1}, \cite{DEK2}. In this paper, we
investigate audio effects \nobreak processing using high-order allpass filters that consist of many cascaded low-order allpass filters. These filters have long chirp-like impulse responses. When audio and music signals are processed with such a filter, remarkable changes are obtained that are similar to the spectral delay effect  \cite{DEK3}, \cite{DEK4}.
\end{extract}
Filtering an audio signal with an allpass filter does not usually have a major effect on the signal's timbre. The allpass filter does not change the frequency content of the signal, but only introduces a phase shift or delay. 
\[
\tau _\textrm{g} (\omega ) =  - \frac{{d\phi (\omega )}}{{d\omega }}.
\]
Audibility of the phase distortion caused by an allpass filter in a sound reproduction system has been a topic of many studies, see, e.g., \cite{DEK1}, \cite{DEK2}. In this paper, we investigate audio effects processing using high-order allpass filters that consist of many cascaded low-order allpass filters. These filters have long chirp-like impulse responses. When audio and music signals are processed with such a filter, remarkable changes are obtained that are similar to the spectral delay effect.
\begin{alphalist}
\item{}Green--function determined experimentally and published.
\item{}Black--function determined using similarity searches and published.
\item{}Red--function determined using similarity searches and determined in this study.
\item{}Blue--O-antigen structure unknown. Function determined using similarity searches and proposed in this study.
\end{alphalist}
Filtering an audio signal with an allpass filter does not usually have a major effect on the signal's timbre. The allpass filter does not change the frequency content of the signal, but only introduces a phase shift or delay. Audibility of the phase distortion caused by an allpass filter in a sound reproduction system has been a topic of many studies, see, e.g., \cite{DEK1}, \cite{DEK2}. 
%Enunciations
\begin{example}
In this paper, we investigate audio effects processing using high-order allpass filters that consist of many cascaded low-order allpass filters. These filters have long chirp-like impulse responses. 
\end{example}
Filtering an audio signal with an allpass filter does not usually have a major effect on the signal's timbre. The allpass filter does not change the frequency content of the signal, but only introduces a phase shift or delay. Audibility of the phase distortion caused by an allpass filter in a sound reproduction system has been a topic of many studies.
\begin{bulletlist}
\item{}Green--function determined experimentally and published.
\item{}Black--function determined using similarity searches and published.
\item{}Red--function determined using similarity searches and determined in this study.
\item{}Blue--O-antigen structure unknown. Function determined using similarity searches and proposed in this study.
\end{bulletlist}
Filtering an audio signal with an allpass filter does not usually have a major effect on the signal's timbre. The allpass filter does not change the frequency content of the signal, but only introduces a phase shift or delay. Audibility of the phase distortion caused by an allpass filter in a sound reproduction system has been a topic of many studies, see, e.g., \cite{DEK1}, \cite{DEK2}. In this paper, we investigate audio effects processing using high-order allpass filters that consist of many cascaded low-order allpass filters. These filters have long chirp-like impulse responses. When audio and music signals are processed with such a filter, remarkable changes are obtained that are similar to the spectral delay effect  \cite{DEK3}, \cite{DEK4}.
\begin{unnumlist}
\item{}Green--function determined experimentally and published.
\item{}Black--function determined using similarity searches and published.
\item{}Red--function determined using similarity searches and determined in this study.
\item{}Blue--O-antigen structure unknown. Function determined using similarity searches and proposed in this study.
\end{unnumlist}
Filtering an audio signal with an allpass filter does not usually have a major effect on the signal's timbre. The allpass filter does not change the frequency content of the signal, but only introduces a phase shift or delay. Audibility of the phase distortion caused by an allpass filter in a sound reproduction system has been a topic of many studies, see, e.g., \cite{DEK1}, \cite{DEK2}. In this paper, we investigate audio effects processing using high-order allpass filters that consist of many cascaded low-order allpass filters. These filters have long chirp-like impulse responses. When audio and music signals are processed with such a filter, remarkable changes are obtained that are similar to the spectral delay effect  \cite{DEK3}, \cite{DEK4}.

\section{SUMMARY}
Filtering an audio signal with an allpass filter does not usually have a major effect on the signal's timbre. The allpass filter does not change the frequency content of the signal, but only introduces a phase shift or delay. Audibility of the phase distortion caused by an allpass filter in a sound reproduction system has been a topic of many studies, see, e.g., \cite{DEK1}, \cite{DEK2}. In this paper, we investigate audio effects processing using high-order allpass filters that consist of many cascaded low-order allpass filters. These filters have long chirp-like impulse responses. When audio and music signals are processed with such a filter, remarkable changes are obtained that are similar to the spectral delay effect  \cite{DEK3}, \cite{DEK4}.

\section{CONCLUSION}
Filtering an audio signal with an allpass filter does not usually have a major effect on the signal's timbre. The allpass filter does not change the frequency content of the signal, but only introduces a phase shift or delay. Audibility of the phase distortion caused by an allpass filter in a sound reproduction system has been a topic of many studies, see, e.g., \cite{DEK1}, \cite{DEK2}. In this paper, we investigate audio effects processing using high-order allpass filters that consist of many cascaded low-order allpass filters. These filters have long chirp-like impulse responses. When audio and music signals are processed with such a filter, remarkable changes are obtained that are similar to the spectral delay effect  \cite{DEK3}, \cite{DEK4}.

\section{ACKNOWLEDGMENT}
This research was conducted in fall 2008 when Vesa V\"alim\"aki was a visiting scholar at CCRMA, Stanford University. His visit was financed by the Academy of Finland (project no. 126310). The authors would like to Dr. Henri Penttinen for his comments and for the snare drum sample used in this work.


Bibliography
\begin{thebibliography}{99}

\bibitem{DEK1}
Preis D., ``Phase Distortion and Phase Equalization in Audio Signal Processing---A Tutorial Review,'' \textit{J. Audio Eng. Soc}, vol. 30, no. 11, pp. 774--794 (1982 Nov.).


\bibitem{DEK2}
Abel J. S. and Berners D. P., ``MUS424/EE367D: Signal Processing Techniques for Digital Audio Effects," Unpublished Course Notes, CCRMA, Stanford University, Stanford, CA (2005).

\bibitem{DEK3}
Roads C., ``Musical Sound Transformation by Convolution," in \textit{Proc. Int. Computer Music Conf.} (Tokyo, Japan, 1993), pp. 102--109.

\bibitem{DEK4}
Roads C., \textit{The Computer Music Tutorial}. MIT Press, Cambridge, MA, 1996.

\end{thebibliography}

%Appendix
\appendix
\section*{APPENDIX}
Filtering an audio signal with an allpass filter does not usually have a major effect on the signal's timbre. The allpass filter does not change the frequency content of the signal, but only introduces a phase shift or delay. Audibility of the phase distortion caused by an allpass filter in a sound reproduction system has been a topic of many studies, see, e.g., \cite{DEK1}, \cite{DEK2}.
\begin{equation}
\phi (\omega ) =  - \omega  + 2\arctan \left( {\frac{{a_1 \sin \omega }}{{1 + a_1 \cos \omega }}} \right)
\end{equation}

In this paper, we investigate audio effects processing using high-order allpass filters that consist of many cascaded low-order allpass filters. These filters have long chirp-like impulse responses. When audio and music signals are processed with such a filter, remarkable changes are obtained that are similar to the spectral delay effect  \cite{DEK3}, \cite{DEK4}.


\begin{nomenclature}[PAMPs]
\subsection*{NOMENCLATURE}
\nomentry{a$_c$}{condensation coefficient condensation coefficient condensation coefficient}


\nomentry{TLR}{Toll-like receptor}

\nomentry{PAMPs}{pathogen-associated molecular patterns condensation coefficient condensation}
\end{nomenclature}

%Biography
 \biography{A1firstname A1lastname}{a.eps}{A1firstname A1lastname is professor of audio signal processing at Helsinki University of Technology (TKK), Espoo, Finland. He received his Master of Science in Technology, Licentiate of Science in Technology, and Doctor of Science in Technology degrees in electrical engineering from TKK in 1992, 1994, and 1995, respectively. His doctoral dissertation dealt with fractional delay filters and physical modeling of musical wind instruments. Since 1990, he has worked mostly at TKK with the exception of a few periods. In 1996 he spent six months as a postdoctoral research fellow at the University of Westminster, London, UK. In 2001-2002 he was professor of signal processing at the Pori School of Technology and Economics, Tampere University of Technology, Pori, Finland. During the academic year 2008-2009 he has been on sabbatical and has spent several months as a visiting scholar at the Center for Computer Research in Music and Acoustics (CCRMA), Stanford University, Stanford, CA. His research interests include musical signal processing, digital filter design, and acoustics of musical instruments. Prof. V\"alim\"aki is a senior member of the IEEE Signal Processing Society and is a member of the AES, the Acoustical Society of Finland, and the Finnish Musicological Society. He was the chairman of the 11th International Conference on Digital Audio Effects (DAFx-08), which was held in Espoo, Finland, in 2008.}
 \biography{A2firstname A2lastname}{b.eps}{A2firstname A2lastname is a consulting professor at the Center for Computer Research in Music and Acoustics (CCRMA) in the Music Department at Stanford University where his research interests include audio and music applications of signal and array processing, parameter estimation, and acoustics. From 1999 to 2007, Abel was a co-founder and chief technology officer of the Grammy Award-winning Universal Audio, Inc. He was a researcher at NASA/Ames Research Center, exploring topics in room acoustics and spatial hearing on a grant through the San Jose State University Foundation. Abel was also chief scientist of Crystal River Engineering, Inc., where he developed their positional audio technology, and a lecturer in the Department of Electrical Engineering at Yale University. As an industry consultant, Abel has worked with Apple, FDNY, LSI Logic, NRL, SAIC and Sennheiser, on projects in professional audio, GPS, medical imaging, passive sonar and fire department resource allocation. He holds Ph.D. and M.S. degrees from Stanford University, and an S.B. from MIT, all in electrical engineering. Abel is a Fellow of the Audio Engineering Society.}
\end{document}